% Part: lambda-calculus
% Chapter: lambda-definability
% Section: primitive-recursive-functions

\documentclass[../../../include/open-logic-section]{subfiles}

\begin{document}

\olfileid[hi]{lam}{ldf}{prf}
\olsection{आदिम पुनरावर्ती फलन \usetoken{S}{lambda definable} हैं}

स्मरण करें कि आदिम पुनरावर्ती फलन वे हैं जिन्हें मूल फलनों $\Zero$,
$\Succ$ और $\Proj{n}{i}$ से संयोजन और आदिम पुनरावर्तन द्वारा परिभाषित किया
जा सकता है।

\begin{lem}
  \ollabel{lem:basic}
  मूल आदिम पुनरावर्ती फलन $\Zero$, $\Succ$ और प्रक्षेप~$\Proj{n}{i}$
  !!{lambda definable} हैं।
\end{lem}

\begin{proof}
उन्हें निम्न पदों द्वारा !!{lambda define} किया जाता है:
\begin{align*}
  \fn{Zero} & \ident \lambd[a][\lambd[fx][x]]\\
  \fn{Succ} & \ident \lambd[a][\lambd[fx][f (a f x)]] \\
  \fn{Proj}^n_i & \ident \lambd[x_0\dots x_{n-1}][x_i]
\end{align*}
\end{proof}

\begin{lem}
  \ollabel{lem:comp} मान लें $k$-स्थानी फलन $f$ और $n$-स्थानी फलन $g_0$,
  \dots, $g_{k-1}$ पदों $F$, $G_0$, \dots, $G_k$ द्वारा !!{lambda
  definable} हैं और $h$ को उनसे संयोजन द्वारा परिभाषित किया गया है। तब $H$
  !!{lambda definable} है।
\end{lem}

\begin{proof}
  ~$h$ को इस पद द्वारा !!{lambda define} किया जा सकता है:
  \[
  H \ident \lambd[x_0 \dots x_{n-1}][F\, (G_0 x_0 \dots
    x_{n-1}) \dots (G_{k-1} x_0 \dots x_{n-1})]
  \]
  इस तथ्य का सत्यापन अभ्यास के लिए छोड़ते हैं।
\end{proof}

\begin{prob}
  यह दिखाकर \olref[lam][ldf][prf]{lem:comp} का प्रमाण पूरा करें कि
  $H\num{n_0}\dots\num{n_{n-1}} \red \num{h(n_0, \dots,
    n_{n-1})}$।
\end{prob}

ध्यान दें कि \olref{lem:comp} ने $f$ और $g_0$, \dots,~$g_{k-1}$ के आदिम
पुनरावर्ती होने की माँग नहीं की; केवल उनका संपूर्ण और !!{lambda definable}
होना आवश्यक है।

\begin{lem}\ollabel{lem:prim}
  मान लें $f$ कोई $n$-स्थानी फलन और~$g$ कोई $n+2$-स्थानी फलन है, वे पदों
  $F$ और~$G$ द्वारा !!{lambda definable} हैं, और फलन~$h$ को $f$ और $g$ से
  आदिम पुनरावर्तन द्वारा परिभाषित किया गया है। तब $h$ भी !!{lambda
  definable} है।
\end{lem}

\begin{proof}
  स्मरण करें कि $h$ इस प्रकार परिभाषित है:
  \begin{align*}
    h(x_1, \dots, x_n, 0) &= f(x_1, \dots, x_n)\\
    h(x_1, \dots, x_n, y+1) & = h(x_1, \dots, x_n, y, h(x_1, \dots, x_n, y)).
  \end{align*}
  अनौपचारिक रूप से, आदिम पुनरावर्ती परिभाषा फलन $h$ के अनुप्रयोग को $y$
  बार पुनरावर्तित करके उसे $f(x_1, \dots, x_n)$ पर लागू करती है। यह चर्च
  अंकों की परिभाषा की याद दिलाता है, जो स्वयं पुनरावर्तक के रूप में
  परिभाषित हैं।

  सरलता के लिए हम एक अतिरिक्त तर्क~$x$ वाली स्थिति की परिभाषा और प्रमाण
  देते हैं। फलन $h$ को इस प्रकार !!{lambda define} किया जाता है:
  \begin{align*}
    H \ident & \lambd[x][\lambd[y][\fn{Snd} (y
        D \tuple{\num{0}, F x})]]
    \intertext{जहाँ}
    D \ident & \lambd[p][\tuple{\fn{Succ} (\fn{Fst}\, p), (G x
          (\fn{Fst}\, p) (\fn{Snd}\, p))}]
  \end{align*}
  हमारे द्वारा रखी गई पुनरावर्तन-अवस्था एक युग्म है, जिसका पहला घटक वर्तमान
  $y$ और दूसरा~$h$ का अनुरूप मान है। पुनरावर्तन के प्रत्येक चरण में हम $y$
  और $h$ के नए मानों का युग्म बनाते हैं; पुनरावर्तन पूरा होने पर युग्म का
  दूसरा भाग लौटाते हैं, जो $h$ का अंतिम मान है। अब हम सिद्ध करते हैं कि यह
  सचमुच आदिम पुनरावर्तन का निरूपण है।

  हम सिद्ध करना चाहते हैं कि किन्हीं $n$ और $m$ के लिए
  $H\,\num{n}\,\num{m} \red \num{h(n,m)}$। इसके लिए पहले दिखाते हैं कि
  यदि $D_n \ident \Subst{D}{\num{n}}{x}$, तो $D_n^m \tuple{\num{0}, F\,
  \num{n}} \red \tuple{\num{m}, \num{h(n, m)}}$। हम~$m$ पर आगमन करते हैं।

  यदि $m=0$, तो हमें $D_n^0 \tuple{\num{0}, F\, \num{n}} \red
  \tuple{\num{0}, \num{h(n, 0)}}$ चाहिए। किंतु $D_n^0
  \tuple{\num{0}, F\, \num{n}}$ केवल $\tuple{\num{0}, F\, \num{n}}$ है।
  चूँकि $F$, ~$f$ को !!{lambda define} करता है, इसलिए यह अपचयित होकर
  $\tuple{\num{0}, \num{f(n)}}$ बनता है; और चूँकि $f(n) = h(n, 0)$, इसलिए
  यह $\tuple{\num{0}, \num{h(n,0)}}$ है।

  अब मान लें $D_n^m \tuple{\num{0}, F\, \num{n}} \red
  \tuple{\num{m}, \num{h(n, m)}}$। हम दिखाना चाहते हैं कि $D_n^{m+1}
  \tuple{\num{0}, F\, \num{n}} \red \tuple{\num{m+1}, \num{h(n,
  m+1)}}$।
  \begin{align*}
    D_n^{m+1} \tuple{\num{0}, F\, \num{n}}
    & \ident D_n(D_n^{m} \tuple{\num{0}, F\, \num{n}})\\
    & \red D_n\,\tuple{\num{m}, \num{h(n, m)}} \text{\qquad (आगमन-परिकल्पना से)}\\
    & \ident (\lambd[p][
      \tuple{
        \fn{Succ} (\fn{Fst}\, p), (G \, \num{n} (\fn{Fst}\, p) (\fn{Snd}\, p))
    }]) \tuple{\num{m}, \num{h(n,m)}}\\
    & \redone
    \tuple{\fn{Succ} (\fn{Fst}\, \tuple{\num{m}, \num{h(n,m)}}),\\
      & \qquad (G \, \num{n} (\fn{Fst}\, \tuple{\num{m}, \num{h(n,m)}}) (\fn{Snd}\, \tuple{\num{m}, \num{h(n,m)}}))}\\
    & \red 
    \tuple{\fn{Succ}\, \num{m}, (G \, \num{n} \,\num{m} \, \num{h(n,m)})}\\
    & \red \tuple{\num{m+1}, \num{g(n, m, h(n, m))}}
  \end{align*}
  चूँकि $g(n, m, h(n, m)) = h(n, m+1)$, प्रमाण पूरा हुआ।

  अंततः, विचार करें:
  \begin{align*}
    H\,\num n\,\num m &\ident \lambd[x][\lambd[y][\fn{Snd} (y
        (\lambd[p].\tuple{\fn{Succ} (\fn{Fst}\, p), (G\, x\,
          (\fn{Fst}\, p)\, (\fn{Snd}\, p))}) \tuple{\num{0}, F x})]]\\
    & \qquad\,\num n\, \num m\\
    & \red \fn{Snd} (\num m \,
    \underbrace{(\lambd[p].\tuple{\fn{Succ} (\fn{Fst}\, p), (G \,\num n\,
      (\fn{Fst}\, p) (\fn{Snd}\, p))})}_{D_n} \tuple{\num{0}, F \num n})\\
    & \ident \fn{Snd} (\num m \, D_n\, \tuple{\num{0}, F \num n})\\
    & \red \fn{Snd} \,(D_n^m  \tuple{\num{0}, F \num n})\\
    & \red \fn{Snd} \,\tuple{\num{m}, \num{h(n, m)}} \\
    & \red \num{h(n,m)}. 
  \end{align*}  
\end{proof}


\begin{prop}
  प्रत्येक आदिम पुनरावर्ती फलन !!{lambda definable} है।
\end{prop}

\begin{proof}
  \olref{lem:basic} से सभी मूल फलन !!{lambda definable} हैं; और
  \olref{lem:comp} तथा \olref{lem:prim} से !!{lambda definable} फलन संयोजन
  और आदिम पुनरावर्तन के अधीन संवृत हैं।
\end{proof}

\end{document}

